% Part I, chapter 2: WFDB records.
\chapter{WFDB records}
\label{chap:records}

A WFDB \emph{record} is a recording of one or more signals, such as the
leads of an ECG, and the files that describe and annotate it: a
\emph{header file}, which describes the record and its signals; one or more
\emph{signal files}, which hold the samples; and any number of
\emph{annotation files}, which label events of the signals, such as beats.
A record may also have an \emph{info file} of comments. This chapter
describes these files as \wfdbx{} reads and writes them, which is as the
WFDB library 10.7.0 does, except where \cref{chap:behavior} says otherwise.
The exact layout of each storage format, and the annotation codes, are in
\cref{chap:formats}.

\section{Records and record names}
\label{sec:record-names}

A record is designated by its \emph{record name}, such as \code{100s}. The
files of a record are found from it:
\begin{itemize}
\item the header is the file \emph{record}\file{.hea} (\file{100s.hea});
\item the signal files are named in the header (\file{100s.dat});
\item an annotation file is the file \emph{record}\code{.}\emph{annotator},
  where the \emph{annotator name} says which set of annotations it holds
  (\file{100s.atr}, the reference annotations of record \code{100s},
  annotator \code{atr});
\item the info file is \emph{record}\file{.info}.
\end{itemize}

A record name may include a directory (\code{mitdb/100},
\code{/data/mitdb/100}), and is then found as \cref{chap:files} describes. A
trailing \file{.hea} is removed from a record name by the library functions
that open a record, so \code{100s.hea} designates record \code{100s}. A name
that ends with \code{/} stands for a record of the same name in that
directory: \code{abc/123/} is \code{abc/123/123}. A record name that ends
with another extension, such as \code{100s.edf}, names an EDF file
(\cref{sec:edf}). The record name \code{-} stands for the standard input or
output (\cref{sec:stdio}).

When a header or an annotation file is created, the record name, and the
annotator name, may contain only letters, digits, \code{\_}, \code{-},
\code{\textasciitilde} and \code{/}; any other character is an error (for
example, \code{init: illegal character 46 in record name}). In a record name
written under a path, such as \code{out.d/rec}, only the last component, the
record's own name, is checked (the C library checked the whole name; see
\cref{sec:resolve}).

\section{The header file}
\label{sec:header}

A header is a text file of lines separated by line feeds; a carriage return
may precede each line feed. Its fields are separated by spaces or tabs. The
lines are of four kinds:
\begin{itemize}
\item the \emph{record line}, the first line that is neither empty nor a
  comment;
\item one \emph{signal specification line} per signal, for an ordinary
  record, or one \emph{segment specification line} per segment, for a
  multi-segment record (\cref{sec:multisegment});
\item \emph{comment lines}, whose first character is \code{\#}; those that
  follow the record line are the \emph{info strings} of the record;
\item empty lines, which are ignored.
\end{itemize}
Comment lines and empty lines may appear anywhere. If the first line of a
header begins with \code{\#wfdb} followed by a version number, such as
\code{\#wfdb 10.7.1}, a library older than that version refuses the header
(\code{init: reading record REC requires WFDB library version 10.7.1 or
later}); \wfdbx{} has the version of the C library that it implements,
10.7.0.

\subsection{The record line}

The record line gives the fields of \cref{tab:record-line}, in this order,
separated by spaces or tabs. Only the record name and the number of signals
are required; each optional field may be given only if the fields before it
are.

\begin{table}[htbp]
\caption{The fields of the record line.}
\label{tab:record-line}
\begin{tabularx}{\linewidth}{@{}P{3.5cm}L@{}}
\toprule
{\tabhead Field} & {\tabhead Meaning} \\
\midrule
record name & The name of the record. It must match the last component of
  the name under which the header is opened, or the header is refused
  (\code{init: record name in record REC header is incorrect}). \\
\code{/}\emph{segments} & (Optional, attached to the record name.) The
  number of segments of a multi-segment record, as in \code{multi/3}. \\
number of signals & An integer; 0 is allowed. Several signals can share one
  signal file. \\
sampling frequency & (Optional.) The number of frames per second, a
  floating-point number greater than 0 (\code{360}, \code{360.0} and
  \code{3.6e2} are equivalent). \wfdbx{} assumes none: if it is missing,
  the frequency is unspecified (0), unless one was set with
  \code{setsampfreq} before the record was opened (the C library assumed
  250\,Hz; see \cref{sec:no-defaults}). \\
\code{/}\emph{counter frequency} & (Optional, attached to the sampling
  frequency.) The frequency of a counter, such as the page numbers of a
  chart recording, used by the time strings that begin with \code{c}
  (\cref{sec:times}). If it is absent or not positive, it is the sampling
  frequency. \\
\code{(}\emph{base counter}\code{)} & (Optional, attached to the counter
  frequency.) The counter value of sample 0; 0 if absent. \\
number of samples & (Optional.) The number of samples per signal. 0 or a
  missing value means that the length is unspecified, and then the
  checksums of the signals are not verified. \\
base time & (Optional.) The time of day of sample 0, as
  \code{hh:mm:ss}, with optional milliseconds (\code{12:34:56.789}). \\
base date & (Optional, after the base time.) The date of sample 0, as
  \code{dd/mm/yyyy}. \\
\bottomrule
\end{tabularx}
\end{table}

\subsection{Signal specification lines}

Each signal specification line gives the fields of
\cref{tab:signal-line}, in this order; again each optional field may be
given only if the fields before it are. The lines of the signals that share
a signal file must be consecutive; such signals form a \emph{signal group},
and must have the same format and block size (otherwise
\code{init: error in specification of signal N or M}).

\begin{table}[htbp]
\caption{The fields of a signal specification line, and the values that
\wfdbx{} uses when a field is missing.}
\label{tab:signal-line}
\begin{tabularx}{\linewidth}{@{}P{3.3cm}L@{}}
\toprule
{\tabhead Field} & {\tabhead Meaning, and value when missing} \\
\midrule
file name & The signal file, relative to the directory of the header
  (\cref{sec:header-files}); \code{-} is the standard input. \\
format & The storage format (\cref{sec:signal-files}), an integer such as
  \code{212}. Required. \\
\code{x}\emph{spf} & (Attached to the format.) The number of samples per
  frame (\cref{sec:frames}); 1 if missing (or less than 1). \\
\code{:}\emph{skew} & (Attached.) The number of samples of the signal that
  precede sample 0, to correct a skew between signals; 0 if missing. \\
\code{+}\emph{offset} & (Attached.) The number of bytes at the beginning of
  the signal file that precede the samples (a prolog); 0 if missing. \\
gain & The number of ADC units per physical unit, a floating-point number.
  0, or a missing field, means that the signal is \emph{uncalibrated}
  (\cref{sec:units}). \\
\code{(}\emph{baseline}\code{)} & (Attached to the gain.) The ADC value
  that corresponds to 0 physical units; the ADC zero if missing. \\
\code{/}\emph{units} & (Attached to the gain or baseline.) The physical
  units, such as \code{mV}, without spaces (at most 50 characters). If
  missing, the units are unspecified; \wfdbx{} never takes them to be
  millivolts, and its tools show them as \code{nd}. \\
ADC resolution & The resolution of the analog-to-digital converter, in
  bits. If missing, it is unspecified (0); the C library inferred it from the
  format or assumed 12 bits. \\
ADC zero & The ADC value at the middle of the ADC's input range; 0 if
  missing. \\
initial value & The value of sample 0, used by the difference format 8;
  the ADC zero if missing. \\
checksum & The sum of all the samples of the signal, as a 16-bit signed
  integer, verified when the end of the signal is read (if the record line
  gives the number of samples); 0 if missing. \\
block size & For signal files that must be read in blocks of a given size
  (character special files); 0 (none) if missing. \\
description & The rest of the line, which may contain spaces (at most 100
  characters), such as \code{MLII}. If missing:
  \code{record REC, signal N}. \\
\bottomrule
\end{tabularx}
\end{table}

\subsection{Info strings and info files}
\label{sec:info-strings}

The \emph{info strings} of a record are the comment lines of its header
that follow the record line, without their initial \code{\#}, as the
library function \code{getinfo} reads them: every line after the record line
whose first character is \code{\#} is an info string, even one between two
signal lines. A line that begins with a space before its \code{\#} is not
an info string. The comment lines before the record line are not info
strings. \code{getinfo} then returns the comment lines of the record's info
file, \emph{record}\file{.info}, if there is one. \code{putinfo} writes an
info string, as \code{\#} followed by the string and a CR LF line ending,
into the header being written, or into the info file opened by
\code{setinfo}. \Cref{chap:metadata} describes the info strings of the form
\code{Key: value}, which hold the metadata of a record.

\subsection{An example: the header of record 100s}

The header of the sample record \code{100s}, \file{tests/data/100s.hea}, is:
\begin{lstlisting}[style=text]
100s 2 360 21600
100s.dat 212 200 11 1024 995 21537 0 MLII
100s.dat 212 200 11 1024 1011 -3962 0 V5
# 69 M 1085 1629 x1
# Aldomet, Inderal
#Produced by xform from record 100, beginning at 0:0
\end{lstlisting}
\begin{itemize}
\item The record line says that record \code{100s} has 2 signals, sampled at
  360 frames per second, of 21600 samples each (60 seconds). It gives no
  base time or date.
\item Each signal line says that the signal is in the file
  \file{100s.dat} (both signals share it: they form signal group 0), in
  format 212, with a gain of 200 ADC units per physical unit, no units, an
  ADC resolution of 11 bits, an ADC zero of 1024, an initial value of 995
  (signal 0) or 1011 (signal 1), a checksum of 21537 or $-3962$, a block
  size of 0, and the description \code{MLII} or \code{V5}. The baseline is
  not given, so it is the ADC zero, 1024. A sample of 995 ADC units in
  signal 0 is thus $(995 - 1024)/200 = -0.145$ physical units.
\item The last three lines are the info strings
  \code{" 69 M 1085 1629 x1"}, \code{" Aldomet, Inderal"} and
  \code{"Produced by xform from record 100, beginning at 0:0"}. The first
  two are free text; the third contains a colon, and is therefore read as a
  metadata entry (\cref{chap:metadata}).
\end{itemize}

\section{Signal files and storage formats}
\label{sec:signal-files}

A signal file holds the samples of the signals of one signal group, as a
sequence of frames (\cref{sec:frames}): the samples of each frame are stored
one signal after the other. The \emph{storage format} of the group says how
each sample is encoded. \Cref{tab:formats-summary} summarizes the formats;
\cref{chap:formats} gives their exact layout.

\begin{table}[htbp]
\caption{The storage formats (summary).}
\label{tab:formats-summary}
\begin{tabularx}{\linewidth}{@{}P{1.4cm}LP{2.6cm}@{}}
\toprule
{\tabhead Format} & {\tabhead Encoding} & {\tabhead Invalid sample} \\
\midrule
0 & no signal file: a null signal, whose samples are all invalid & (all) \\
8 & 8-bit first differences, from the initial value of the header & --- \\
16 & 16-bit two's complement, least significant byte first & $-32768$ \\
24 & 24-bit two's complement, least significant byte first & $-2^{23}$ \\
32 & 32-bit two's complement, least significant byte first & $-2^{31}$ \\
61 & 16-bit two's complement, most significant byte first & $-32768$ \\
80 & 8-bit offset binary (128 is subtracted from each byte) & $-128$ \\
160 & 16-bit offset binary, least significant byte first & $-32768$ \\
212 & 12-bit two's complement, two samples packed in 3 bytes & $-2048$ \\
310 & 10-bit two's complement, three samples packed in 4 bytes (two 16-bit
  words) & $-512$ \\
311 & 10-bit two's complement, three samples packed in a 32-bit word & $-512$ \\
508, 516, 524 & FLAC-compressed, with 8, 16 or 24 bits per sample (at most
  8 signals per file, all with the same number of samples per frame) &
  $-128$, $-32768$, $-2^{23}$ \\
\bottomrule
\end{tabularx}
\end{table}

The last column is the value that marks an invalid (missing) sample in the
file: the most negative value of the format. The library returns an invalid
sample, whatever its format, as \code{WFDB\_INVALID\_SAMPLE} ($-32768$), and
the high-level readers give it as NaN in physical units. A file in format
212, 310 or 311 may end with a partial group of samples. FLAC files are read
and written by \wfdbx{} without libFLAC: the files it writes decode to the
same samples as those of the C library, but their compressed bytes differ.

Three fields of the signal line change how a signal file is read: the
\emph{byte offset} skips a prolog at the beginning of the file; the
\emph{skew} discards the first samples of a signal, to align it with the
others (the skipped samples are counted in the checksum); and the
\emph{block size} makes the file be read in blocks of that size.

\section{ADC units and physical units}
\label{sec:units}

A signal file stores samples as integers, in \emph{ADC units} (adu), the
output of the analog-to-digital converter. The header gives, for each
signal, the values that relate them to \emph{physical units}, such as
millivolts:
\begin{description}
\item[gain] $g$, the number of ADC units per physical unit (0: the signal
  is uncalibrated);
\item[baseline] $b$, the ADC value that corresponds to 0 physical units;
\item[units] the physical units, such as \code{mV} (unspecified if the
  header does not give them);
\item[ADC resolution] $r$, in bits, and \textbf{ADC zero} $z$, the value at
  the middle of the ADC's range. Together they give the range of the ADC:
  from $z - 2^{r-1}$ to $z + 2^{r-1} - 1$.
\end{description}

The conversions are those of the C library, and of the high-level
functions:
\begin{itemize}
\item From ADC units to physical units (\code{aduphys(s, a)}, and
  \code{read\_record} with physical units): $v = (a - b)/g$, as a
  floating-point number. An invalid sample becomes NaN in the high-level
  readers.
\item From physical units to ADC units (\code{physadu(s, v)}, and
  \code{write\_record} with physical units): $a = \mathrm{round}(v g) + b$,
  where $\mathrm{round}(x)$ is $\lfloor x + 0.5\rfloor$ for $x \ge 0$ and
  $\lceil x - 0.5\rceil$ for $x < 0$ (halves are rounded away from zero).
  In \code{write\_record}, NaN becomes an invalid sample.
\item \code{adumuv(s, a)} converts $a$ ADC units to microvolts, as
  $\mathrm{round}(1000\,a/g)$, and \code{muvadu(s, v)} converts $v$
  microvolts to ADC units, as $\mathrm{round}(g\,v/1000)$. Neither uses the
  baseline; both take the physical units to be millivolts.
\end{itemize}
In these functions, $g$ and $b$ are those of input signal \code{s} of the
open record. \wfdbx{} assumes no gain: for an uncalibrated signal (gain 0)
or a signal number that is not that of an input signal, the four functions
report an error (\code{aduphys: signal 0 is uncalibrated},
\code{aduphys: invalid signal number 5}) and return NaN (\code{aduphys}) or
\code{WFDB\_INVALID\_SAMPLE} (the other three); the C library assumed a gain
of 200 (\cref{sec:no-defaults}).

\section{Frames, samples per frame and multi-frequency records}
\label{sec:frames}

A \emph{frame} is the set of samples of all the signals that belong to one
sampling interval of the record. The sampling frequency of the record line
is the frequency of the frames. Usually each signal has one sample per
frame. A signal sampled at an integer multiple $n$ of that frequency has
$n$ samples per frame (the \code{x}\emph{n} suffix of its format), and its
own sampling frequency is $n$ times that of the frames. A record with such
signals is a \emph{multi-frequency} record. For example, the test record
\file{tests/data/mfreq.hea}:
\begin{lstlisting}[style=text]
mfreq 3 100 600
mfreq.dat 16 200/mV 16 0 138 -16869 0 slow
mfreq.dat 16x2 200/mV 16 0 218 5410 0 medium
mfreq.dat 16x4 200/mV 16 0 219 -15996 0 fast
\end{lstlisting}
has 600 frames at 100 frames per second; the signals are sampled at 100,
200 and 400\,Hz, and each frame of its signal file holds 1 + 2 + 4 = 7
samples.

The library reads such a record in one of two modes, set by
\code{setgvmode}:
\begin{itemize}
\item \textbf{Low resolution} (\code{WFDB\_LOWRES}, the default):
  \code{getvec} returns one sample per signal per frame. The samples of a
  signal that has more than one per frame are averaged: their mean is
  rounded to the nearest integer, halves upward, and is invalid if one of
  them is invalid. \code{sampfreq} returns the frame frequency.
\item \textbf{High resolution} (\code{WFDB\_HIGHRES}): \code{getvec}
  returns as many sample vectors per frame as the largest number of samples
  per frame (4 for \code{mfreq}). Each holds the next sample of each signal
  that has that many samples per frame; the samples of the other signals
  are repeated. \code{sampfreq} returns the frame frequency times that
  number (400\,Hz for \code{mfreq}), and times (\cref{sec:times}) count
  these sample vectors.
\end{itemize}
\code{WFDB\_GVPAD}, which can be combined with either mode, replaces each
invalid sample by the previous valid sample of the same signal.
\code{getframe} returns all the samples of a frame. In the high-level API,
\code{read\_record} reads in low resolution by default, and in high
resolution with its \code{highres} option; its \code{pad} option is
\code{WFDB\_GVPAD}. For \code{mfreq}, the first frame is read as
\code{138 1459 499} in low resolution (1459 is the mean of 218 and 2699,
rounded up), and as these four rows in high resolution:
\begin{lstlisting}[style=text]
138   218    219
138   218  -1316
138  2699   1120
138  2699   1973
\end{lstlisting}
The library can also resample all the signals of a record to another
frequency, by linear interpolation (\code{setifreq}).

\section{Multi-segment records}
\label{sec:multisegment}

A multi-segment record is a sequence of ordinary records, its
\emph{segments}, read as if their signals were continuous. Its header (the
\emph{master header}) has a record line whose record name is followed by
\code{/} and the number of segments, and then one segment specification
line per segment: the record name of the segment and its length in samples.
The segments' own headers and signal files are found in the directory of the
master header (\cref{sec:header-files}). A segment cannot itself be a
multi-segment record (\code{init: record REC cannot be nested in another
multi-segment record}), and the sampling frequency must be the same in all
segments. If the record line gives a number of samples that differs from the
sum of the segment lengths, the library warns
(\code{warning (init): in record REC, stated record length (N)} \ldots).

In a \emph{fixed-layout} record, every segment has the same signals, in
the same order, with the same gains and baselines (their storage formats
can differ). The test record \file{tests/data/multi.hea} is one: two copies
of record \code{100s} with a null record between them.
\begin{lstlisting}[style=text]
multi/3 2 360 45000
100s 21600
null 1800
100s 21600
\end{lstlisting}

In a \emph{variable-layout} record, the signals may differ from segment to
segment. Segment 0, of length 0, is then the \emph{layout segment}: its
header, whose signal file names are \code{\textasciitilde}, gives the
signals of the record, with their gains, baselines and units. When the
record is read, the signals of each segment are rearranged, shifted and
rescaled to the gains and baselines of the layout (the results rounded to
the nearest integer, halves upward); a signal absent from a segment reads as
invalid samples. If a segment's gain is
larger than that of the layout, the library warns
(\code{sigmap\_init: loss of precision in signal N in segment SEG}). A
segment whose record name is \code{\textasciitilde} is a \emph{null
segment}: it has no files, and all its samples are invalid. The test record
\file{tests/data/vlay.hea} and its layout segment \file{vlay\_0.hea}:
\begin{lstlisting}[style=text]
vlay/4 3 250 3201
vlay_0 0
fmt16 1500
~ 200
fmt212 1501
\end{lstlisting}
\begin{lstlisting}[style=text]
vlay_0 3 250 0
~ 0 400(10)/mV 12 0 0 0 0 sig0 fmt16
~ 0 400/uV 12 0 0 0 0 sig1 fmt212
~ 0 100/mmHg 16 0 0 0 0 sig2 fmt16
\end{lstlisting}

\section{Annotation files}
\label{sec:annotations}

An annotation file holds a sequence of \emph{annotations}, each of which
labels a point of the record, such as a beat. It is the file
\emph{record}\code{.}\emph{annotator}. Each annotation has the fields of
\cref{tab:annotation-fields}, those of the C structure
\code{WFDB\_Annotation}.

\begin{table}[htbp]
\caption{The fields of an annotation.}
\label{tab:annotation-fields}
\begin{tabularx}{\linewidth}{@{}P{2.2cm}L@{}}
\toprule
{\tabhead Field} & {\tabhead Meaning} \\
\midrule
\code{time} & The time of the annotation, in sample intervals from the
  beginning of the record (a 64-bit integer). \\
\code{anntyp} & The \emph{annotation code}, from 1 to \code{ACMAX} (49),
  such as \code{NORMAL} (1) or \code{PVC} (5). Each code has a
  \emph{mnemonic}, such as \code{N} or \code{V}, and a description
  (\code{annstr}, \code{anndesc}); \cref{chap:formats} lists them. \\
\code{subtyp} & The subtype, a signed 8-bit integer ($-128$ to 127). \\
\code{chan} & The channel (signal) number, an unsigned 8-bit integer (0 to
  255). \\
\code{num} & The annotator number, a signed 8-bit integer. \\
\code{aux} & Optional auxiliary data, up to 255 bytes, such as the rhythm
  \code{(N} of a \code{RHYTHM} annotation, or the text of a \code{NOTE}. In
  the library functions, its first byte is its length, as in C. \\
\bottomrule
\end{tabularx}
\end{table}

The high-level \code{read\_annotations} returns these fields as arrays:
the times (\code{samples}), the codes, their mnemonics (\code{symbols}),
the subtypes, the channels, the annotator numbers and the aux strings.

Annotation files are in one of two binary formats, which the library tells
apart when it opens a file (with a warning, \code{file appears to be in AHA
format}, if the caller asked for the other):
\begin{itemize}
\item \textbf{MIT format}, the usual one: each annotation is a 16-bit word,
  least significant byte first, whose 6 high bits are the code and 10 low
  bits the interval from the previous annotation, followed by optional
  words for the fields that change (pseudo-codes \code{SKIP} for a long
  interval, \code{NUM}, \code{SUB}, \code{CHN} and \code{AUX}); a word of 0
  ends the file.
\item \textbf{AHA format}: 16 bytes per annotation, as in the distribution
  files of the AHA database.
\end{itemize}

An annotation file is in \emph{canonical order} when its annotations are
sorted by time, then annotator number, then channel. An output file that is
not is sorted when it is closed (\cref{sec:sorting}). When a program has set
a time resolution with \code{setafreq} (as \cmd{wrann} does), the library
writes it at the beginning of the file, as a \code{NOTE} annotation whose
aux string is \code{\#\# time resolution: F}. Annotation files can also
begin with \emph{modification labels}, \code{NOTE} annotations at time 0
that redefine the mnemonics and descriptions of codes; the library reads
them when it opens the file, so \code{getann} does not return them.

\section{Times and dates}
\label{sec:times}

Times are counted in \emph{sample intervals} from the beginning of the
record (in high resolution, in intervals of the fastest signal;
\cref{sec:frames}). The library function \code{strtim} converts a time
string to a number of sample intervals, and \code{timstr} and
\code{mstimstr} convert back; the tools take their time arguments
(\opt{-f}, \opt{-t}, \ldots) as \code{strtim} strings, and the high-level
readers accept them for \code{sampfrom} and \code{sampto}.

\subsection{Time strings}

\code{strtim} first skips spaces, tabs, line feeds and carriage returns, and
then reads the string according to its first character, as in
\cref{tab:strtim}. In the table, $f$ is the sampling frequency (the one set
by \code{setifreq}, if any), and a number is rounded to an integer by adding
0.5 and truncating toward zero.

\begin{table}[htbp]
\caption{The time strings that \code{strtim} accepts.}
\label{tab:strtim}
\begin{tabularx}{\linewidth}{@{}P{3.6cm}L@{}}
\toprule
{\tabhead String} & {\tabhead Value (sample intervals)} \\
\midrule
\code{s}\emph{n}, as \code{s123} & Sample number $n$. \\
\code{e} & The end of the record: its number of samples (times the number
  of samples per frame of the fastest signal, in high resolution); 0 if it
  is not known. \\
\code{i} & The time of the next sample to be read from the input
  signals. \\
\code{o} & The time of the next sample to be written to the output
  signals. \\
\code{c}\emph{x}, as \code{c350.5} & Counter value $x$: $(x - x_0)\,f/f_c$,
  with the counter frequency $f_c$ and base counter $x_0$ of the header.
  With no counter frequency, $x$ is read as a sample number. \\
\code{f}\emph{n}, as \code{f100} & Frame number $n$: $n f/f_0$, where $f_0$
  is the frame frequency (in high resolution, the number of the first sample
  of that frame). \\
\code{[}\emph{time}\code{]}, as \code{[13:06:00]} & A time of day: the
  value is the \emph{negated} number of sample intervals from the base time
  of the record to that time, so that it is 0 or negative. \\
\code{[}\emph{time} \emph{d}\code{]}, as \code{[8:00:00 1]} & A time of day
  $d$ days after the base date. \\
\code{[}\emph{time} \emph{date}\code{]}, as \code{[12:00:00 1/3/1992]} & A
  time of day on a date (\code{dd/mm/yyyy}). \\
\emph{s}, as \code{143} or \code{14.875} & An elapsed time of $s$ seconds:
  $s f$. \\
\emph{m}\code{:}\emph{s}, as \code{2:14.875} & $m$ minutes and $s$
  seconds. \\
\emph{h}\code{:}\emph{m}\code{:}\emph{s}, as \code{4:02:01} & $h$ hours, $m$
  minutes and $s$ seconds. \\
\bottomrule
\end{tabularx}
\end{table}

A string that cannot be read gives 0, which is also a valid result
(\code{strtim("0")}). The high-level readers make the value of a time of day
positive, so that it is the number of the sample at that time. For record
\code{100s} (360\,Hz), \code{"0:10"} and \code{"10"} are 3600, \code{"s123"}
is 123, \code{"1:02:03.5"} is 1340460, \code{"e"} is 21600, and
\code{"[00:00:10]"} is $-3600$.

\wfdbx{} assumes no sampling frequency: when none is known (no record is
open, or its header gives none, and none was set with \code{setsampfreq}),
\code{strtim} reports \code{strtim: sampling frequency is not specified} and
returns 0 for every string that needs one, that is, for all the strings of
\cref{tab:strtim} except those that begin with \code{s}, \code{e},
\code{i} or \code{o}, and those that begin with \code{c} when there is no
counter frequency. The function \code{wfdbx\_time\_needs\_freq(string)}
returns 1 if \code{strtim} needs the frequency for a string, and 0
otherwise. The C library used 1 sample per second instead
(\cref{sec:no-defaults}).

\subsection{Formatted times}

\code{timstr(t)} converts a time $t$ to text:
\begin{itemize}
\item if $t > 0$, or if the record has no base time or date, as an elapsed
  time $|t|/f$, truncated to whole seconds: \code{h:mm:ss}, or, under one
  hour, \code{mm:ss}, right-aligned in 8 characters (\code{"    0:10"} for
  10 seconds);
\item if $t \le 0$ and the record has a base time or date, as the time of
  day at $-t$ sample intervals after the base time, in brackets:
  \code{[hh:mm:ss dd/mm/yyyy]}, with the date if the record has a base date
  (or the number of days after the base time if it has none).
\end{itemize}
\code{mstimstr(t)} does the same with milliseconds (\code{"    0:10.003"},
\code{[12:34:56.789 31/12/1999]}). Without a sampling frequency, both report
\code{timstr: sampling frequency is not specified} (or \code{mstimstr: \ldots})
and return the time as a sample number, \code{s} followed by $|t|$, which
\code{strtim} reads back.

\subsection{Base time and dates}

The base time and date of a record come from its header, or from
\code{setbasetime}, which takes a string \code{hh:mm:ss}, with optional
milliseconds and an optional date \code{dd/mm/yyyy} after a space (an empty
string sets the current local date and time). \code{strdat} and
\code{datstr} convert between dates \code{dd/mm/yyyy} and day numbers. The
high-level readers give the base time as text (\code{base\_time}, as
\code{timstr(0)} gives it, to the second), as a number of milliseconds after
midnight (\code{base\_time\_ms}), and the base date as year, month and day
(\code{base\_date}).

\section{Calibration files}
\label{sec:calibration}

A calibration file lists, for types of signals, the calibration pulses with
which they were recorded and their customary plotting scale. Each line that
is not a comment (\code{\#}) has six fields:
\begin{lstlisting}[style=text]
DESC<tab>LOW HIGH TYPE SCALE UNITS
\end{lstlisting}
\code{DESC} is a signal type, matched against the beginning of a signal's
description (\code{*} matches any signal); \code{LOW} and \code{HIGH} are
the levels of the calibration pulse in physical units (\code{-} for
\code{LOW}: the signal is AC-coupled; \code{-} for \code{HIGH}: not
defined); \code{TYPE} is \code{sine}, \code{square}, \code{sawtooth} or
\code{undefined}; \code{SCALE} is in physical units per centimeter; and
\code{UNITS} are the physical units. For example:
\begin{lstlisting}[style=text]
ECG	- 1 sine 1 mV
BP	0 100 square 100 mmHg
\end{lstlisting}
\code{calopen(name)} reads a calibration file into the library's
calibration list (emptied first, unless the name begins with \code{+});
without a name, it reads the file named by the environment variable
\code{WFDBCAL}, or \file{wfdbcal}. The file is found as any other file
(\cref{chap:files}). \code{getcal} finds the first entry for a signal's
description and units, \code{putcal} adds one, and \code{newcal} writes the
list to a file. The file \file{tests/data/wfdbcal} is the default
calibration file of the C package.

\section{EDF and EDF+ records}
\label{sec:edf}

The library reads EDF and EDF+ files (and BDF files, with 24-bit samples)
directly, as records. A record name with an extension other than
\file{.hea}, such as \code{100s.edf}, is opened as an EDF file: its header
is read from the EDF header, and its samples from the data records that
follow. Each EDF data record is one frame, so a signal has as many samples
per frame as it has samples per data record, and the frame frequency is the
inverse of the duration of a data record. The signals are in format 16 (24
for BDF). The gain of a signal is
$(d_{\max} - d_{\min})/(p_{\max} - p_{\min})$, from its digital and physical
ranges (0 if the physical range is empty), its baseline is
$d_{\max} - p_{\max} \times \mathrm{gain}$ rounded to the nearest integer,
and its ADC zero is the middle of the digital range. The record's base time
and date are the start time and date of the EDF file.

Because frames are data records, the library opens an EDF file in high
resolution (\cref{sec:frames}): the sample intervals are those of the
fastest signal, and \code{sampfreq} gives its frequency. A program that sets
another mode reads one sample per data record per signal; the high-level
\code{read\_record}, which sets the mode, returns every sample of an EDF file
only with its \code{highres} option.

The command-line tools \cmd{edf2mit} and \cmd{mit2edf} convert EDF files to
WFDB records and back, and \cmd{rdedfann} prints the annotations of an EDF+
file in the format of \cmd{rdann}.
